\textbf{Phần 1. Trắc nghiệm lựa chọn}
\setcounterpageref{ex}{0}

\begin{ex} %[3D11N62]
Số gồm $7$ chục nghìn, $2$ nghìn, $0$ trăm và $5$ đơn vị là
\choice
{$72050$}
{$72500$}
{\True $72005$}
{$70205$}
\loigiai{
Số đó là $72005$.\\
\textbf{Đáp án của câu hỏi là C}
}
\end{ex}

\begin{ex} %[3D11N62]
Số liền trước của số $100\,000$ là
\choice
{$99990$}
{$90000$}
{$99900$}
{\True $99999$}
\loigiai{
$100000 - 1 = 99999$.\\
\textbf{Đáp án của câu hỏi là D}
}
\end{ex}

\begin{ex} %[3D11H62]
Điền dấu thích hợp: $56\,700 \dots 57\,600$.
\choice
{\True $<$}
{$>$}
{$=$}
{Không so sánh được}
\loigiai{
Hàng nghìn: $6 < 7$ nên $56700 < 57600$.\\
\textbf{Đáp án của câu hỏi là A}
}
\end{ex}

\begin{ex} %[3D11H62]
Làm tròn số $48\,750$ đến hàng nghìn ta được
\choice
{$48000$}
{\True $49000$}
{$50000$}
{$48700$}
\loigiai{
Hàng trăm là 7 nên làm tròn lên $49000$.\\
\textbf{Đáp án của câu hỏi là B}
}
\end{ex}

\begin{ex} %[3D11H62]
Làm tròn số $73\,200$ đến hàng chục nghìn ta được
\choice
{\True $70000$}
{$80000$}
{$73000$}
{$75000$}
\loigiai{
Hàng nghìn là 3 nên làm tròn thành $70000$.\\
\textbf{Đáp án của câu hỏi là A}
}
\end{ex}

\begin{ex} %[3D11H62]
Số lớn nhất trong các số $84\,500; 85\,400; 84\,050; 85\,040$ là
\choice
{$84500$}
{\True $85400$}
{$85040$}
{$84050$}
\loigiai{
Số lớn nhất là $85400$.\\
\textbf{Đáp án của câu hỏi là B}
}
\end{ex}

\begin{ex} %[3D11V62]
Từ các chữ số $4, 0, 9, 2, 5$ lập số bé nhất có 5 chữ số khác nhau là
\choice
{\True $20459$}
{$24059$}
{$20495$}
{$02459$}
\loigiai{
Số bé nhất là $20459$.\\
\textbf{Đáp án của câu hỏi là A}
}
\end{ex}

\begin{ex} %[3D11V62]
Dãy số nào sau đây xếp theo thứ tự tăng dần?
\choice
{$45000; 43000; 42000; 41000$}
{\True $32000; 34000; 36000; 38000$}
{$89000; 87000; 85000; 83000$}
{$50000; 40000; 30000; 20000$}
\loigiai{
Dãy tăng dần: $32000 < 34000 < 36000 < 38000$.\\
\textbf{Đáp án của câu hỏi là B}
}
\end{ex}

\begin{ex} %[3D11V62]
Tìm chữ số $x$ để: $7x\,500 > 78\,500$.
\choice
{$8$}
{$7$}
{\True $9$}
{$6$}
\loigiai{
$x > 8 \Rightarrow x = 9$.\\
\textbf{Đáp án của câu hỏi là C}
}
\end{ex}

\begin{ex} %[3D11V62]
Một thư viện có $45\,680$ cuốn sách. Làm tròn số sách đến hàng nghìn ta được khoảng
\choice
{$45000$ cuốn}
{$50000$ cuốn}
{$45700$ cuốn}
{\True $46000$ cuốn}
\loigiai{
Hàng trăm là 6 nên làm tròn lên $46000$ cuốn.\\
\textbf{Đáp án của câu hỏi là D}
}
\end{ex}

\begin{ex} %[3D11V62]
Số liền sau của số $79\,999$ là
\choice
{$79998$}
{$80001$}
{\True $80000$}
{$89999$}
\loigiai{
$79999 + 1 = 80000$.\\
\textbf{Đáp án của câu hỏi là C}
}
\end{ex}

\begin{ex} %[3D11C62]
Tìm số tròn chục nghìn có 5 chữ số bé nhất mà tổng các chữ số của nó bằng 8.
\choice
{$80000$}
{$17000$}
{$26000$}
{\True $80000$ (vì chữ số hàng chục nghìn là $8$)}
\loigiai{
Số tròn chục nghìn có dạng $\overline{a0000}$. Tổng các chữ số bằng 8 nên $a = 8 \Rightarrow 80000$.\\
\textbf{Đáp án của câu hỏi là D}
}
\end{ex}

\textbf{Phần 2. Trắc nghiệm đúng sai}
\setcounterpageref{ex}{0}

\begin{ex} %[3D11H62]
Xét tính đúng sai về số học trong phạm vi 100 000:
\choiceTF
{\True Số $100\,000$ là số tròn trăm nghìn.}
{\True $65\,400 < 65\,450$.}
{Làm tròn $45\,000$ đến hàng chục nghìn được $40\,000$.}
{\True Làm tròn $45\,000$ đến hàng chục nghìn được $50\,000$.}
\loigiai{
\begin{itemize}
    \item $100000$ tròn trăm nghìn: Đúng.
    \item $65400 < 65450$: Đúng.
    \item Hàng nghìn là 5 nên làm tròn thành 50000. Ý 3 Sai, Ý 4 Đúng.
\end{itemize}
\begin{itemize}
    \item a) \textbf{Đúng}.
    \item b) \textbf{Đúng}.
    \item c) \textbf{Sai}.
    \item d) \textbf{Đúng}.
\end{itemize}
}
\end{ex}

\begin{ex} %[3D11V62]
Cho 3 số: $54\,800; 58\,400; 54\,080$. Xét tính đúng sai:
\choiceTF
{\True Số lớn nhất là $58\,400$.}
{\True Số bé nhất là $54\,080$.}
{Làm tròn số $54\,800$ đến hàng nghìn được $54\,000$.}
{\True Làm tròn số $54\,800$ đến hàng nghìn được $55\,000$.}
\loigiai{
\begin{itemize}
    \item Lớn nhất: $58400$. Đúng.
    \item Bé nhất: $54080$. Đúng.
    \item Hàng trăm là 8 nên làm tròn thành 55000. Ý 3 Sai, Ý 4 Đúng.
\end{itemize}
\begin{itemize}
    \item a) \textbf{Đúng}.
    \item b) \textbf{Đúng}.
    \item c) \textbf{Sai}.
    \item d) \textbf{Đúng}.
\end{itemize}
}
\end{ex}

\textbf{Phần 3. Trả lời ngắn}
\setcounterpageref{ex}{0}

\begin{ex} %[3D11H62]
Viết số liền trước của số $90\,000$.
\shortans{89999}
\loigiai{
\textbf{Đáp số:} $89999$\\
$90000 - 1 = 89999$.
}
\end{ex}

\begin{ex} %[3D11H62]
Viết số liền sau của số $89\,999$.
\shortans{90000}
\loigiai{
\textbf{Đáp số:} $90000$\\
$89999 + 1 = 90000$.
}
\end{ex}

\begin{ex} %[3D11V62]
Làm tròn số $64\,780$ đến hàng nghìn được số nào?
\shortans{65000}
\loigiai{
\textbf{Đáp số:} $65000$\\
$65000$.
}
\end{ex}

\begin{ex} %[3D11V62]
Làm tròn số $75\,200$ đến hàng chục nghìn được số nào?
\shortans{80000}
\loigiai{
\textbf{Đáp số:} $80000$\\
$80000$.
}
\end{ex}

\begin{ex} %[3D11V62]
Tìm số lớn nhất có 5 chữ số khác nhau lập từ các chữ số $8, 1, 5, 0, 9$.
\shortans{98510}
\loigiai{
\textbf{Đáp số:} $98510$\\
$98510$.
}
\end{ex}

\begin{ex} %[3D11C62]
Tìm số tròn chục nghìn $x$ biết $45\,000 < x < 65\,000$.
\shortans{50000}
\loigiai{
\textbf{Đáp số:} $50000$\\
$50000$ (hoặc $60000$).
}
\end{ex}
